\section{问题分析}

\subsection{总体思路}

四个问题共享同一套物理对象与同一组物理约束，差别只在两处：\textbf{决策时刻能获得哪些信息}，以及\textbf{计划与实际不一致时如何结算}。因此本文采用“\emph{一个物理内核 $+$ 四条信息集主线}”的建模框架：先把母线能量平衡、储能状态转移、容量与功率限值、弃光上界写成公共约束组，再对每一问替换相应费用项、扩展决策变量，并重新界定信息集。这样既保持四问口径一致，也能把结果差异归因到信息更新与结算规则本身。

具体地，问题 1 是完全信息下的确定性经济调度，用于验证物理约束与求解器；问题 2 在\textbf{完全无法获得当天预报}的极端信息集下，仅用历史观测构造因果预测，并把 $5$ 倍紧急电价转化为非对称损失；问题 3 引入官方多时刻预报，形成“$0{:}00$ 基准计划 $\rightarrow$ 每 $6$ h 滚动调整 $\rightarrow$ 实际执行 $\rightarrow$ 状态递推”的闭环；问题 4 只把电价序列换成逐日逐时的波动电价，用于检验前两类策略在价格波动下的稳健性。

整体流程为“数据对齐与单位换算 $\rightarrow$ 公共物理模型 $\rightarrow$ 因果预测与滚动重优化 $\rightarrow$ 主口径求解 $\rightarrow$ 消融、回代与提交文件核验”，各环节共享同一数据索引和费用定义。

\subsection{具体分析}

\indent\textbf{问题一：}题目给定附件 1 的电价、负载与光伏发电预测，且“每天的电价和小区负载相同”，因此属于确定性优化。富余光伏必须经储能充电或弃光消纳，故公共模型保留弃光变量；本问的最优解实际实现零弃光。题目还要求 $0{:}00$ 与 $24{:}00$ 储电量相同，该闭环条件必须作为硬约束显式写入。

\indent\textbf{问题二：}题目未提供任何当天预报，本文主模型按最严格信息边界，仅依赖已经发生的历史观测构造因果预测。若业务另行提供日前全天负荷，则可采用负荷完全已知的口径；本文将其作为理想信息基准，用于量化两种信息条件的费用差异。$5$ 倍电价使“高估光伏”的局部边际损失是“低估光伏”的 $4$ 倍，据此可在单时段报童损失下给出保守分位点；全年调度则采用因果点预测的确定性等价模型。计划在 $0{:}00$ 锁定后，实际偏差由储能进行 $10$ min 级实时缓冲。

\indent\textbf{问题三：}本问引入 $6{:}00$、$12{:}00$、$18{:}00$ 三个新增预报时刻，题目据此询问“是否需要引入”。本文把“允许更新决策的节点集合”作为消融变量，构造 $S_0$（仅 $0{:}00$）$\to S_1$（$+6{:}00$）$\to S_2$（$+12{:}00$）$\to S_3$（$+18{:}00$）四组对照；节点未启用时，最近一次有效计划持续执行，绝不回退到 $0{:}00$ 基准计划。由于启用节点会同时获得已实现负载观测与新发布光伏预报，主消融量度的是\textbf{整个决策更新节点的价值}；再用配对控制实验分别识别负载刷新与光伏预报刷新的贡献。题目将计划费、调整费与紧急购电费并列相加，故叠加结算作为主口径，差额结算仅作敏感性对照，两者分别重新优化。

\indent\textbf{问题四：}本问按题面字面要求，直接把附件 4 给定的逐日实时电价代入问题二、三重新计算，并将这一口径作为主答案与结果工作簿来源。附件 4 的电价在日内与日间均有明显波动（最低 $0.0076$ 元/kWh，最高 $1.7936$ 元/kWh），峰谷价差扩大了储能套利空间。考虑到实际市场中完整日内价格的发布时间可能受交易规则限制，本文另设只使用已发生价格的因果预测情形作为现实性扩展；二者之差用于说明价格预知信息的价值，但不替代题面主结果。

\subsection{求解框架}

问题一、二及差额结算对照为线性规划；问题三和问题四的叠加结算主模型加入充放电互斥的 $0$--$1$ 变量，形成混合整数线性规划。二者均由 HiGHS 求解。连续模型再以“真实费用—储能吞吐量—弃光量”词典序目标排除无经济意义的退化最优解；叠加模型由互斥约束直接禁止同一时段充放电。所有电量按 $10$ min 时段能量（kWh）核算，功率与能量换算系数为 $\Delta t=1/6$ h。
